\documentclass[12pt]{article}
\usepackage[utf8]{inputenc}
\usepackage[T1]{fontenc}
\usepackage{hyperref}
\title{Touo Language Resources for K-12 Curriculum Development}
\author{Compiled by Claude Sonnet 5 (claude-sonnet-5) for Andr\'as Kornai}
\date{}
\begin{document}
\maketitle

\section{Introduction}
Touo (ISO 639-3 tqu; Glottocode touo1238), also called Baniata or Lokuru, is spoken on southern Rendova Island, Western Province, Solomon Islands. Speaker estimates in the sources found are about 1,900 (1999 census) and 1,870. The task description gives about 2,800. Its affiliation is disputed: it is often placed in a tentative Central Solomons family, but Glottolog treats it as an isolate. Speakers are typically multilingual in other languages, Solomon Islands Pijin and English. One source states that schools function in English and there is no institutional support for Touo literacy.

\section{Primary Educational Resources}
No textbooks or primers were found. The only item is a study of vernacular literacy:
\begin{itemize}
\item Vernacular literacy in the Touo language of the Solomon Islands (Max Planck Institute for Psycholinguistics publication record): \url{https://www.mpi.nl/publications/item_58044}
\end{itemize}

\section{Linguistic and Reference Materials}
\begin{itemize}
\item Touo language (Wikipedia): \url{https://en.wikipedia.org/wiki/Touo_language}
\item Omniglot page on Touo: \url{https://omniglot.com/writing/touo.htm}
\item Touo numerals (Max Planck Institute for Evolutionary Anthropology): \url{https://lingweb.eva.mpg.de/channumerals/Touo.htm}
\end{itemize}

\section{Religious and Community Resources}
None found.

\section{Digital and Audio Resources}
\begin{itemize}
\item Digital documentation project preserving Touo oral narratives (University of Colorado Boulder, Linguistics): \url{https://www.colorado.edu/linguistics/2023/04/07/digital-documentation-preserving-touo-oral-narratives-across-globe}
\end{itemize}

\section{Curriculum Development Context}
Touo is stable and universally used as a first language, but schooling is in English. The Colorado documentation project and the literacy study are the most useful starting points. They suggest a community-based orthography and story-based reading materials. A usable dictionary, grammar or primer was not found in these searches.
\end{document}